\section{Pray, Fight, Work}

The historiography of the Right needs to be presented in two parts. Evola has written about one of them, specifically, the need to understand historical and current events as manifestations of the degeneration of caste. The other, and more important, is to show how a spiritual elite was able to lead a civilization in a Traditional direction. Here, a thorough knowledge of the Medieval civilization is of value, since we can use actual texts and historical events, rather than rely on reconstructions from myths and legends.

To use this historiography, it is necessary to keep in mind that caste is the most fundamental typology of humans, prior even to race and ethnicity. Evola, himself, makes that clear, since his ``international right'' would transcend ethnic unities. Although, because of the confusion and undifferentiation of modernity, one's own caste may be difficult to discern. Nevertheless, it is an essential part of one's nature (not meant in a biological sense).

To do that task, I will be relying on the book \textit{The Three Orders}, by \textbf{Georges Duby}, an historian in the lineage of \textbf{George Dumezil}. Dumezil\footnote{\url{https://gornahoor.net/index.php?s=dumezil}}, in his lifelong study of Indo-European civilizations from India to Iceland and Ireland, demonstrated what he called trifunctionality. That is, such civilizations were characterized by a division into the three functions: those who pray, those who fight, and those who work. Obviously, one way trifunctionality is manifested is through a caste system. Now Rene Guenon claimed that caste divisions\footnote{\url{https://gornahoor.net/?p=1698}} are a touchstone of Traditional civilizations, and thus has metaphysical significance.

By limiting himself to a period of two centuries in one area of France, Duby is able to show how this idea of trifunctionality was understood by the spiritual authorities and the policies they implemented to organize their society along traditional lines. We will also use \textbf{Ananda Coomaraswamy}'s book Spiritual Authority and Temporal Power in the Indian Theory of Government to show the correspondence between the Medieval and Vedic conceptions of the three functions.

\textbf{Gerard of Cambrai} and \textbf{Adalbero of Laon} are the two bishops that Duby focuses on, in the period of about the year 1000 to 1200. Thomas Bisson's introduction makes clear what Duby is doing:

\begin{quotationx}
This is a book about the political and cultural uses of a social idea. It comes on as a doubly original book: first, because it appears to be the first book ever devoted to \emph{the history of what became the paradigmatic image of the ancien regime}; and second, because it seeks to explain the appearance and early diffusion of that image as expressions of the strategies of threatened or \emph{innovating elites}.

The tripartite conception of society is one of those collection ``imaginings'' of which the records should be read not only in the light of historical actuality but also to reveal those structural articulations of human experience, with their continuities and interruptions, \emph{which inform a cultural history running, in this case, from Indo-European antiquity to the French Revolution}.
\end{quotationx}


Of course, Duby, writing as an historian, can do no more than describe ideas and events; whether or not the bishops were expounding true metaphysical principles. Nevertheless, even as a merely descriptive principle, Duby demonstrates the remarkable continuity of the Traditional nature of European peoples from their known origins right up to the French Revolution; he has no religious or metaphysical position to defend. Immediately, we see that those pseudo-traditionalists who deny the Traditional character of the Medieval period are way off the mark. Not only do we have the writings of the major intellects of Tradition, but now we have the historical facts themselves.

We see also that those who defend the ancien régime, are defending the memory and values of the European people, for nearly the totality of their history. We will see what well-bred men of that time considered to be healthy and normal. We will see the relationships between the castes and the intellectual justification for its structure. There may be wise men today who will learn from this study how to duplicate the accomplishments of the bishops in their own time.


\flrightit{Posted on 2012-11-07 by Cologero}

\begin{center}* * *\end{center}

\begin{footnotesize}
\begin{sffamily}
\texttt{Logres on 2012-11-14 at 21:57 said:}

Adolbert hailed from Reims, where Hincmar had objected to the papacy that they were trying to enslave the free Franks, a generation earlier. Judging by his name, he was not Gallo-Roman, and Adolbert objected to persons of humble birth becoming bishops. From reading Maurras, it would appear that the Franks were attempting to convert their entire society to a ``feudalistic'' basis en masse: that is, a total society organized towards God. This is consistent with Charlemagne's views on conversions, when he forcibly baptized the Saxons. It was (in this sense) anti-Christian by modern standards, but rooted in Tradition. That is, it believed in authority being born (Maurras) when the Franks stepped in to uphold law and order in the wake of Roman collapse. It would appear that the West Romans desired to tame this, as later the papacy would likewise try to tame the Germanic emperors.


\hfill

\end{sffamily}
\end{footnotesize}

